% ============================================================
%  preamble.tex —— 全项目共用导言区（main.tex 与分片测试均 \input 本文件）
%  在 \documentclass[11pt,a4paper,UTF8,fontset=fandol]{ctexart} 之后使用
% ============================================================

% ---------- 页面与字体 ----------
\usepackage[a4paper,top=2.6cm,bottom=2.6cm,left=2.6cm,right=2.6cm]{geometry}
\usepackage{amsmath,amssymb,amsfonts,amsthm}
\usepackage{graphicx}
\usepackage{float}
\usepackage{caption}
\captionsetup{font=small,labelsep=quad}
\usepackage{booktabs}
\usepackage{longtable}
\usepackage{array}
\usepackage{multirow}
\usepackage{enumitem}
\usepackage{xcolor}
\usepackage{tcolorbox}
\tcbuselibrary{breakable}
\usepackage{fancyhdr}
\usepackage[colorlinks=true,linkcolor=blue!40!black,citecolor=blue!40!black,urlcolor=blue!40!black,
            bookmarksnumbered=true,bookmarksopen=true]{hyperref}
\hypersetup{
  pdftitle={弦紧致化的几何工具（TASI 2017 讲座）——中文译注本},
  pdfauthor={Lara B. Anderson, Mohsen Karkheiran（原著）；中文译注本},
  pdfsubject={弦紧致化、Calabi--Yau 流形、全纯向量丛、模空间},
  pdfkeywords={string compactification, Calabi-Yau, heterotic string, vector bundles, moduli, TASI 2017}
}

% ---------- 版式 ----------
\setlength{\parskip}{0.35ex}
\setlength{\emergencystretch}{3em}
\setlist{itemsep=0.2ex,topsep=0.5ex}
\pagestyle{fancy}
\fancyhf{}
\fancyhead[L]{\small 弦紧致化的几何工具（TASI 2017 讲座）}
\fancyhead[R]{\small\thepage}
\renewcommand{\headrulewidth}{0.4pt}

% ---------- 本节重点（Key Takeaways）----------
\newtcolorbox{keytakeaway}{
  title={本节重点（Key Takeaways）},
  fonttitle=\small\bfseries,
  coltitle=black,
  colbacktitle=gray!18,
  colback=gray!6,
  colframe=gray!45,
  boxrule=0.5pt,
  arc=1.5pt,
  left=6pt,right=6pt,top=4pt,bottom=4pt,
  fontupper=\small,
  breakable
}
\newcommand{\takeawaylabel}{\textbf{本节重点}\quad}

% ---------- 图表路径 ----------
\graphicspath{{figures/}}

% ---------- 常用符号（不改变原文记号，仅便于录入）----------
\newcommand{\Oh}{\mathcal{O}}
\newcommand{\bbZ}{\mathbb{Z}}
\newcommand{\bbC}{\mathbb{C}}
\newcommand{\bbP}{\mathbb{P}}
\newcommand{\bbR}{\mathbb{R}}
\newcommand{\bpa}{\bar\partial}
